
\documentclass[11pt,a4paper]{article}

\usepackage[T1]{fontenc}
\usepackage[utf8]{inputenc}

\usepackage[a4paper,margin=2.7cm]{geometry}
\usepackage{microtype}
\usepackage{mathtools,amssymb,amsthm}

\title{On a Congruence Involving the Decimal Digit Sum}
\author{Matteo Calvigioni}
\date{}

\newcommand{\N}{\mathbb{N}}

\theoremstyle{definition}
\newtheorem{definition}{Definition}

\begin{document}

\maketitle

\begin{abstract}
We consider the congruence
\[
2S(n) - p \equiv n \pmod{p},
\]
where \(S(n)\) denotes the sum of the decimal digits of \(n\).
This relation is not a universal identity; rather, for each fixed modulus \(p\),
it defines a subset of integers satisfying the congruence.
\end{abstract}

\section{Introduction}

Let \(S(n)\) denote the sum of the decimal digits of a positive integer \(n\).
For a fixed integer \(p \ge 2\), we study the congruence
\[
2S(n) - p \equiv n \pmod{p}.
\]
Since \(p \equiv 0 \pmod{p}\), this is equivalent to
\[
2S(n) \equiv n \pmod{p}.
\]
Thus, the problem reduces to comparing the digit sum \(S(n)\) with \(n\) modulo \(p\).

\section{A classical fact on the digit sum}

Let \(n\) be a positive integer. We write its decimal (base-10) positional representation as
\[
n = \sum_{i=0}^{k} a_i 10^i, \qquad a_i \in \{0,1,\dots,9\}.
\]
We define the digit sum of \(n\) as
\[
S(n) = \sum_{i=0}^{k} a_i.
\]

A classical property is that
\[
n \equiv S(n) \pmod{9}.
\]

Indeed, we can write
\[
n - S(n) = \sum_{i=0}^{k} a_i (10^i - 1),
\]
and since \(10 \equiv 1 \pmod{9}\), it follows that \(10^i \equiv 1 \pmod{9}\) for all \(i\), hence each term is divisible by 9, and therefore
\[
n \equiv S(n) \pmod{9}.
\]

\section{Example: the case \(p = 3\)}

Since \(3 \mid 9\), we have
\[
S(n) \equiv n \pmod{3}.
\]
Substituting this into
\[
2S(n) \equiv n \pmod{3},
\]
we obtain
\[
2n \equiv n \pmod{3},
\]
and therefore
\[
n \equiv 0 \pmod{3}.
\]

So, in the case \(p = 3\), the original congruence holds exactly for the multiples of \(3\):
\[
2S(n) - 3 \equiv n \pmod{3}
\quad \Longleftrightarrow \quad
3 \mid n.
\]

\section{Active integers}

\begin{definition}
For a fixed modulus \(p\), define
\[
A_p \coloneqq \{\, n \in \N : 2S(n) - p \equiv n \pmod{p} \,\}.
\]
We call the elements of \(A_p\) the \emph{active integers} with respect to \(p\).
\end{definition}

\section{Discussion}

The congruence
\[
2S(n) - p \equiv n \pmod{p}
\]
should therefore be viewed as a selection condition rather than as a general identity.

Its behavior depends strongly on the modulus \(p\):
\begin{itemize}
    \item for \(p = 3\), the active integers are precisely the multiples of \(3\);
    \item for moduli related to \(9\), the digit sum naturally interacts with congruences;
    \item for general \(p\), the structure of \(A_p\) appears less regular and may be explored computationally.
\end{itemize}

\section{Conclusion}

For each fixed modulus \(p\), the condition
\[
2S(n) - p \equiv n \pmod{p}
\]
defines a subset \(A_p \subseteq \N\).
The example \(p=3\) shows that this subset can sometimes be described explicitly.
A natural continuation is to study the structure of \(A_p\) for other values of \(p\).

\end{document}
